\section{\bt{Design of \system}{\system\ 的设计}}
\label{sec:model}
\bi{\system's design connects two sources of savings: agents reuse one another's discoveries, and the layer shares the work of keeping those discoveries usable. Following Figure~\ref{fig:architecture}, we explain how definitions are learned and published, what makes their aggregation correct, and how validation, repair, and optimization are shared. We then extend the same representation to predictive features with per-instance time bounds.}
{\system\ 的设计把两类节省连接起来：智能体复用彼此的发现，共享层则分摊维持这些发现可用的工作。沿图~\ref{fig:architecture}，本节先说明定义如何学习和发布、其聚合何时正确，再介绍验证、修复与优化如何共享，最后将同一表示扩展到带逐样本时间边界的预测特征。}

\subsection{\bt{What agents share and how it is learned}{智能体共享什么、如何学来}}
\label{sec:shared-object}
\bi{\emph{Shared objects.} \system\ shares database knowledge that agents discover through exploration and whose structural assumptions can be checked against data: metric definitions, including their grain, joins, filters, and encoded value meanings. It also shares table descriptions, verified join paths, and check results to avoid repeated database work. Question parameters and answers remain separate; a result cache handles exact query repeats. Agent strategies, conversations, and preferences remain in agent memory.}
{\emph{共享对象。}\system\ 共享由智能体探索发现、其结构前提可由数据检查的数据库知识：指标定义，包括粒度、连接、过滤及其编码的取值含义。它也共享表描述、验证过的连接路径与检查结果，避免重复的数据库工作。题目参数与答案不属于定义；结果缓存处理完全相同的查询。智能体策略、对话与偏好仍保存在智能体记忆中。}

\bi{A definition revision is $m^r=(B,I,C,E,r)$, where $B$ is its business meaning, $I$ its structured implementation, $C$ its set of executable conditions, $E$ its learning record, and $r$ its revision number. The implementation records the fact table, aggregate expression, grain, date join, joins, and definition filters; its joins form a many-to-one tree rooted at the fact table (a star or snowflake). Question parameters such as the month are not part of it: the generated SQL is compiled from $I$ for any period, difference, or ranking, so one learning session serves questions that no agent has asked yet. Figure~\ref{fig:metric-record} illustrates how this record supports reuse and maintenance.}
{定义的一个版本记为 $m^r=(B,I,C,E,r)$，其中 $B$ 为业务含义，$I$ 为结构化实现，$C$ 为可执行条件集，$E$ 为学习记录，$r$ 为版本号。实现记录事实表、聚合表达式、粒度、日期连接、连接和定义过滤；连接构成以事实表为根的多对一连接树（星型或雪花型）。月份等题目参数不属于实现：生成的 SQL 由 $I$ 针对任意期间、差值或排名编译得到，因此一次学习会话可以服务尚无智能体问过的问题。图~\ref{fig:metric-record} 展示该记录如何支持复用与维护。}

\bi{\emph{Extraction and publication.} Publication turns a successful answer into a definition that other agents can execute. The session must state the business definition and have an independently verified answer, supplied by a user, report, or analyst. In our benchmarks, an automated oracle verifies the final scalar answer against the reference on the same snapshot before extraction. The LLM extractor then uses the question, SQL and arithmetic trace, rejected queries, and schema and join metadata to propose a JSON record; business rows and result values are withheld. The server derives evidence and conditions from its fields and applies seven deterministic publication gates (G1--G7) in three stages.}
{\emph{提取与发布。}发布将一次成功作答转化为其他智能体可以执行的定义。会话须明确业务定义，并有由用户、报表或分析师提供的独立核实答案。在基准实验中，自动判题器先将最终标量答案与同一快照上的参考答案比较，再开始提取。LLM 提取器根据题目、SQL 与算术轨迹、被拒绝的查询以及模式和连接元数据提出 JSON 记录，不接收业务行或结果值。服务器从记录字段推导证据与条件，分三个阶段执行七项确定性发布门控（G1--G7）。}

\bi{\textbf{\emph{(1) Provenance checks.}} G1 requires an explicit business-definition basis; G2 requires an independently judged answer and a verified calculation chain.}
{\textbf{\emph{（1）来源验证。}}G1 要求明确的业务定义依据；G2 要求答案经过独立判定，且计算链可验证。}

\bi{\textbf{\emph{(2) Structural and grain checks.}} G3 verifies valid tables, columns, joins, expressions, and definition filters, with question parameters excluded. G4 checks fact-grain uniqueness. G5 requires a single read-only example query that passes execution review.}
{\textbf{\emph{（2）结构与粒度检查。}}G3 验证表、列、连接、表达式与定义过滤合法，并排除题目参数。G4 检查事实粒度唯一性。G5 要求示例是通过执行审查的单条只读查询。}

\bi{\textbf{\emph{(3) Replay checks.}} G6 confirms that learning dependencies are unchanged and that the example reproduces the verified value. G7 confirms value reproduction by SQL compiled from the structured fields, testing the artifact consumers will execute and catching omissions introduced during extraction. Numeric comparisons use the learning task's declared precision.}
{\textbf{\emph{（3）执行重放。}}G6 确认学习依赖未变化，且示例重现核实值。G7 确认由结构化字段编译的 SQL 重现该值，直接检验使用者将执行的工件，发现提取时遗漏的条件。数值比较使用学习任务声明的精度。}

\bi{\emph{Use.} An agent queries \code{find\_metric} with a metric name or phrase, matched against canonical names and aliases extracted during learning. The agent receives the valid revision: its business definition, structured fields, generated SQL, and revision number. It sets its own period in the generated SQL, or wraps the SQL for a difference or a ranking, declares the revision to \code{run\_sql}, and answers. This replaces the exploration: the agent needs no table listing, value probe, or join test. On TPC-DS, such a session takes three LLM calls for a metric whose definition needs three inner joins; agents without memory spend \PpNoneMTwoSteps\ calls on questions about it, because they must find the fact table, the measure, and every join themselves. Agents issue at most three memory lookups per question before falling back to standard schema exploration.}
{\emph{使用。}智能体向 \code{find\_metric} 提交指标名或短语，与学习过程中提取的规范名称及别名匹配。智能体得到有效版本：业务定义、结构化字段、生成的 SQL 与版本号。它在生成的 SQL 中填入自己的期间，或在外面包一层求差值或排名，在 \code{run\_sql} 中声明该版本，然后作答。这就取代了探索：智能体不必列表、探查取值或检验连接。在 TPC-DS 上，对一个定义需要三个内连接的指标，这样的会话只用三次大模型调用；无记忆的智能体在关于它的题上要花 \PpNoneMTwoSteps\ 次调用，因为事实表、度量和每个连接都要自己找。智能体每道题最多查找三次记忆，未命中时转入常规模式探索。}

\subsection{\bt{When shared knowledge is correct}{共享的知识何时正确}}
\label{sec:sufficiency}
\bi{Data-only updates can break learned SQL while the schema stays the same. Status-change rows, restatements, and duplicate loads violate the grain; reused dimension keys make a many-to-one join fan out; late or reformatted date keys move facts to another period or to none~\cite{kimball2013toolkit}. A condition $c$ is a predicate over the tables $\operatorname{dep}(c)$ with a check query $Q_c$ that returns violating rows or nothing. Its normalized form $\operatorname{id}(c)$ consists of its type, tables, keys, join direction, normalized filters, and null policy. \system\ derives the following four condition types from $I$.}
{纯数据更新可在模式不变的情况下破坏学到的 SQL。状态流水、更正保留旧行或重复加载会破坏粒度；被重用的维度键使多对一连接放大行数；迟到或改了格式的日期键把事实移到别的期间或使其无处可归~\cite{kimball2013toolkit}。条件 $c$ 是表集合 $\operatorname{dep}(c)$ 上的谓词，附带一个检查查询 $Q_c$，返回违例行或空。其规范化形式 $\operatorname{id}(c)$ 由类型、表、键、连接方向、规范化的过滤和空值规则组成。\system\ 从 $I$ 推导以下四类条件。}

\bi{\textbf{\emph{(1) Filtered key uniqueness.}} No two rows of $\sigma_\varphi(T)$ may agree on $K$, as checked by filtered key uniqueness $\operatorname{KeyUnique}(T,K,\varphi)$.}
{\textbf{\emph{（1）过滤后键唯一性。}}过滤后键唯一性 $\operatorname{KeyUnique}(T,K,\varphi)$ 要求 $\sigma_\varphi(T)$ 中没有两行在 $K$ 上取值相同。}

\bi{\textbf{\emph{(2) Join cardinality.}} This condition bounds the number of $R$ rows matched by each row of $L$ in a join $L\to R$.}
{\textbf{\emph{（2）连接基数。}}该条件限制连接 $L\to R$ 中每条 $L$ 记录匹配的 $R$ 记录数。}

\bi{\textbf{\emph{(3) Date join.}} The fact's date key must join the date dimension many-to-one.}
{\textbf{\emph{（3）日期连接。}}事实的日期键须与日期维度多对一连接。}

\bi{\textbf{\emph{(4) Join completeness.}} For the date join and every inner join, the fraction of fact rows that match no row must stay within the fraction recorded at publication plus 0.1 percentage points.}
{\textbf{\emph{（4）连接完整性。}}对日期连接和每个内连接，匹配不到任何行的事实行比例不得超过发布时记录的比例加 0.1 个百分点。}

\bi{For a period $p$, the generated SQL evaluates $\mathrm{AGG}(e)$ over the rows of $\sigma_\varphi(F)\bowtie R_1\bowtie\cdots\bowtie R_k$ whose date falls in $p$. The business definition establishes two domain premises connecting physical records to business events: (B1) the metric's business events correspond one-to-one to the values of $K$ in $\sigma_\varphi(F)$, and (B2) $e$ evaluates on an event's row to that event's measurement.}
{对期间 $p$，生成的 SQL 在 $\sigma_\varphi(F)\bowtie R_1\bowtie\cdots\bowtie R_k$ 中日期落在 $p$ 内的行上计算 $\mathrm{AGG}(e)$。业务定义确立两条连接物理记录与业务事件的领域前提：（B1）指标的业务事件与 $\sigma_\varphi(F)$ 中 $K$ 的取值一一对应；（B2）在事件对应的行上，$e$ 求值得到该事件的度量值。}
\begin{proposition}[\bt{Sufficiency}{充分性}]
\label{prop:sufficiency}
\bi{If B1 and B2 hold on $D_s$ and every condition in $C(m^r)$ holds on $D_s$, then for every period $p$ the generated SQL of $m^r$ returns $\mathrm{AGG}$ over the measurements of the events whose date lies in $p$, each event contributing exactly once; the events it omits are those whose date or join key matches no row, and the completeness conditions bound their fraction.}
{若 B1、B2 在 $D_s$ 上成立，且 $C(m^r)$ 中每个条件都在 $D_s$ 上成立，则对每个期间 $p$，$m^r$ 的生成的 SQL 返回日期落在 $p$ 内的事件度量值的 $\mathrm{AGG}$，每个事件恰好贡献一次；被遗漏的只有日期键或连接键匹配不到任何行的事件，其比例受完整性条件约束。}
\end{proposition}
\begin{proof}[\bt{Proof sketch}{证明概要}]
\bi{By B1 and filtered key uniqueness, $\sigma_\varphi(F)$ holds one row per event. Join cardinality keeps every join from duplicating a row, join completeness bounds the rows that inner joins drop, and the date join gives each row at most one date and hence at most one period. Every remaining row contributes its measurement (B2) once, to the period of its date.\ifbilingual\else\qedhere\fi}
{由 B1 和过滤后键唯一性，$\sigma_\varphi(F)$ 中每个事件恰有一行。连接基数保证连接不复制任何行，连接完整性约束内连接排除的行，日期连接给每行至多一个日期，因而至多一个期间。于是每个保留下来的行把其度量值（B2）恰好贡献一次，归入其日期所在的期间。\qedhere}
\end{proof}
\bi{Conversely, while B1 and B2 hold, every data-only update that makes the generated SQL count an event twice or drop facts beyond the recorded fraction violates a condition.}
{反过来，只要 B1、B2 成立，凡是使生成的 SQL 把一个事件计数两次、或排除超过记录比例的事实的纯数据更新，都会违反某个条件。}

\subsection{\bt{Maintenance once for all agents}{为所有智能体一次性维护}}
\label{sec:maintenance}
\bi{\emph{Shared check results.} Definitions that rely on the same structural condition can share its validation. We identify a condition by $\operatorname{id}(c)$ and the state it reads by the versions of its dependency tables. A table version $V(T)$ combines a schema fingerprint, DML counters, and an ETL batch identifier, or uses a \emph{transactional version} incremented by a statement-level trigger inside each writing transaction. Keying the result by $\bigl(\operatorname{id}(c),V(\operatorname{dep}(c))\bigr)$ lets every definition and agent reuse it at those versions.}
{\emph{共享检查结果。}依赖同一结构条件的定义可以共享验证。我们用 $\operatorname{id}(c)$ 标识条件，用依赖表的版本标识检查所读的状态。表版本 $V(T)$ 可以由结构指纹、DML 计数与 ETL 批次号组成，也可以使用\emph{事务性版本}，由语句级触发器在每个写事务内递增。将检查结果以 $\bigl(\operatorname{id}(c),V(\operatorname{dep}(c))\bigr)$ 为键保存后，所有定义和智能体都能在这些版本上复用它。}

\bi{On lookup, \system\ retains results for unchanged dependencies and searches for completed results at the new versions of changed dependencies. It runs $Q_c$ only when a result is missing, coalescing concurrent requests into that execution. Shared join entries receive the same treatment, so a changed or vanished join path makes the definition pending. A failed condition invalidates the revision; consumers then explore until a successful learning session publishes it again.}
{查找定义时，\system\ 保留依赖未变的检查结果，并为变化的依赖查找新版本上的已完成结果。只有缺少结果时才执行 $Q_c$，并将并发请求合并到同次执行。共享连接条目也采用这一过程，因此连接路径变化或消失会使定义进入待验证状态。条件失败则使版本失效，使用者转回探索，直到成功的学习会话重新发布定义。}

\bi{\emph{Update cost and scheduling.} In the transactional-version implementation, a statement-level trigger performs one metadata upsert per modified table and writing statement, distributed over 16 counter shards to reduce write contention. Validation runs lazily on demand: after an update, the first consumer triggers and completes any uncached dependency checks before use; concurrent consumers coalesce onto that execution, and subsequent consumers reuse its results while dependency versions match.}
{\emph{更新代价与调度。}事务性版本实现中，语句级触发器对每个被修改的表、每条写语句执行一次元数据 upsert，分散到 16 个计数分片以减少写争用。验证采用按需惰性执行：更新后，首个消费者在使用前触发并完成未缓存的依赖检查；并发消费者合并到同次执行，后续消费者在依赖版本一致时复用其结果。}

\bi{\emph{Repair once.} A fact table can lose its grain because an update adds superseded or duplicate rows. Deterministic repair searches equality filters on low-cardinality columns for a way to restore one row per business key without losing keys. A candidate is accepted when exactly one filter passes and the repaired definition reproduces the learning query's result through an \emph{as-of query} on the learning snapshot. Requiring one passing filter distinguishes an unambiguous repair from competing primary and backup copies. The repair serves every agent, and requests arriving during it join the same execution. If no repair can be accepted, the definition remains invalid until a session relearns it.}
{\emph{只修复一次。}更新加入被更正的旧行或重复行时，事实表可能失去原有粒度。确定性修复在低基数列上搜索等值过滤，尝试恢复每个业务键一行，同时保留所有键。候选需要满足两个验收条件：恰好一个过滤通过，并且修复后的定义通过学习快照上的 \emph{as-of 查询}重现学习查询的结果。只接受一个通过的过滤，是为了将明确的修复与主数据、备份副本之间的歧义区分开。一次修复服务所有智能体，期间到来的请求合并到同次执行；若没有可接受的修复，定义保持失效，直到后续会话重新学到它。}

\bi{\emph{Checks at query execution.}\label{sec:execution} Validation must describe the data a consumer actually reads. A writer can commit between an earlier check and query execution, so \system\ validates the declared revisions inside the query's repeatable-read transaction. It looks up condition results at the snapshot's transactional versions, runs missing checks there, and executes the query if they hold. \code{run\_sql} also rejects absent or invalid revisions and SQL missing a required grain filter. Reusing existing results keeps this snapshot check inexpensive; the following lemma establishes when that reuse is valid.}
{\emph{执行时检查。}验证需要描述使用者实际读取的数据。早先的检查与查询执行之间可能有写入提交，因此 \system\ 在查询自己的可重复读事务中验证所声明的版本。它按快照中的事务性版本查找条件结果，在该事务内补做缺失的检查，条件成立后执行查询。\code{run\_sql} 同时拒绝不存在或已失效的版本，以及缺少必需粒度过滤的 SQL。复用已有结果使快照检查保持较低开销；下面的引理说明这种复用何时成立。}

\label{sec:snapshot}
\begin{lemma}[\bt{Check-result reuse}{检查结果复用}]
\label{lem:reuse}
\bi{For read-only snapshots under the write-coverage assumptions below, let $v_s(T)$ be the transactional version of $T$ visible in snapshot $s$. If $v_{s_1}(T)=v_{s_2}(T)$ for every $T\in\operatorname{dep}(c)$, then $c(D_{s_1})=c(D_{s_2})$.}
{对于满足下述写入覆盖前提的只读快照，设 $v_s(T)$ 为快照 $s$ 中可见的 $T$ 的事务性版本。若对每个 $T\in\operatorname{dep}(c)$ 都有 $v_{s_1}(T)=v_{s_2}(T)$，则 $c(D_{s_1})=c(D_{s_2})$。}
\end{lemma}
\begin{proof}
\bi{The visible committed-transaction sets are nested in snapshot acquisition order, say $s_1\subseteq s_2$. A transaction that writes $T$ increments $T$'s counter in the same transaction, so equal versions leave no transaction in $s_2\setminus s_1$ that writes $T$, and both snapshots see the same state of every table that $Q_c$ reads.}
{按快照获取顺序，可见的已提交事务集合相互包含，不妨设 $s_1\subseteq s_2$。写 $T$ 的事务在同一事务内递增 $T$ 的计数，因此版本相等意味着 $s_2\setminus s_1$ 中没有写 $T$ 的事务，两个快照看到 $Q_c$ 所读各表的同一状态。}
\end{proof}
% PostgreSQL 18: sql-copy.html (triggers/FREEZE), sql-createtrigger.html
% (foreign-key actions), sql-truncate.html (transactional but not MVCC-safe).
\bi{\emph{Practical write coverage in PostgreSQL.} On the instrumented base tables used here, without partitioning or inheritance, the trigger covers inserts, updates, deletes (including \code{COPY FROM} and cascading foreign-key actions), and \code{TRUNCATE}; counter increments commit or roll back with the write. Reuse requires counters never to be reset or reused, so committed $A\!\to\!B\!\to\!A$ changes advance the version. Two integration rules complete this coverage: operations outside ordinary MVCC visibility (\code{TRUNCATE}, table rewrites, and \code{COPY FREEZE}) require quiescent readers; DDL, trigger-bypassing writes or replication, and restore/reset operations require cache invalidation before reuse. Schema fingerprints and ETL batch markers supplement detection, while reuse guarantees rely on observing every write.}
{\emph{PostgreSQL 中的实际写入覆盖。}对本文使用的已安装触发器且无分区或继承的基表，插入、更新、删除（包括 \code{COPY FROM} 与外键级联操作）及 \code{TRUNCATE} 均触发版本递增，计数与写入一起提交或回滚。复用要求计数不重置、不复用，因此已提交的 $A\!\to\!B\!\to\!A$ 变化仍推进版本。两条集成规则补全覆盖：绕过普通 MVCC 可见性的操作（\code{TRUNCATE}、表重写及 \code{COPY FREEZE}）要求没有在途读事务；DDL、绕过触发器的写入或复制，以及恢复／重置操作，要求在复用前清除缓存检查结果。模式指纹与 ETL 批次标记辅助检测，复用保证仍以观测到每次写入为前提。}

\subsection{\bt{Optimization once for all agents}{为所有智能体一次性优化}}
\label{sec:optimization}
\bi{Optimizing a shared definition improves the SQL retrieved by every consumer. After learning, \system\ searches the valid definitions for two rewrites: \emph{key range} replaces a period join with a range on the fact table's date key, whose bounds scalar subqueries obtain from the date dimension; \emph{left-join removal} drops a left join that supplies no column and no filter. Candidates retain the structured definition form, so the compiler, conditions, and Proposition~\ref{prop:sufficiency} continue to apply.}
{优化共享定义，可以改善每个使用者检索到的 SQL。学习后，\system\ 在有效定义中搜索两种改写：\emph{日期键范围}用事实表日期键上的范围谓词替换期间连接，上下界由标量子查询从日期维度获取；\emph{去掉左连接}移除不提供列也不带过滤的左连接。候选保留结构化定义的形式，因此编译器、条件和命题~\ref{prop:sufficiency} 仍然适用。}

\bi{The optimizer evaluates rule candidates first. If none is adopted, an optional LLM stage proposes candidates from the definition, generated SQL, execution plan, and table sizes. Publication accepts only structures generated by the two rules, whose equivalence is established by Proposition~\ref{prop:rewrite}; other proposals remain suggestions. Each candidate passes the static, grain, and SQL checks before its assumptions, results, and cost are tested.}
{优化器先评估规则候选；若未采纳任何候选，可选的 LLM 阶段再根据定义、生成的 SQL、执行计划与表规模提出候选。发布只接受两条规则生成的结构，其等价性由命题~\ref{prop:rewrite} 确立；其他提议保留为建议。每个候选先通过静态、粒度和 SQL 检查，再检验前提、结果与代价。}

\bi{For key range, the added assumption is date-key contiguity, checked before adoption and maintained afterwards. Candidate and current revision must agree for the learning question on the learning snapshot and, within one repeatable-read snapshot, on the learning year's twelve months, a quarter, a period difference, and a monthly ranking. Cost is measured in that snapshot over at least eight periods, alternating execution order. Adoption requires a 95\% confidence interval for the paired time difference below zero and a mean saving of at least 10\%.}
{日期键范围改写新增日期键连续性前提，采纳前检查，之后持续维护。候选与当前版本须在学习快照上对学习题返回相同结果，并在同一可重复读快照内，对学习年份的十二个月、一个季度、一次期间差值和一次月度排名返回相同结果。代价也在该快照中测量，覆盖至少八个期间，交替执行候选和当前版本。只有配对耗时差的 95\% 置信区间低于零、平均节省至少 10\% 时，候选才被采纳。}
\begin{proposition}[\bt{Rule equivalence}{规则等价}]
\label{prop:rewrite}
\bi{Let the date dimension's key be unique on $D$. (a) If \emph{date-key contiguity} holds on $D$, i.e., the keys of every month form an integer interval that begins one above the previous month's last key, then for every period of consecutive months the generated SQL of a key-range revision returns on $D$ what that of its previous revision returns. (b) If the right side of a removed left join is unique on the join key on $D$, removing the join does not change the result.}
{设日期维度的键在 $D$ 上唯一。（a）若\emph{日期键连续性}在 $D$ 上成立，即每个月的键构成一个整数区间、且从上个月最后一个键加一开始，则对由连续月份组成的每个期间，键区间版本的生成的 SQL 在 $D$ 上返回与其前一版本相同的结果。（b）若被消除的左连接右侧在连接键上唯一，消除该连接不改变结果。}
\end{proposition}
\begin{proof}[\bt{Proof sketch}{证明概要}]
\bi{(a) By contiguity, the keys of a period's months are exactly the integer interval between the bounds, so a fact row's key lies in the interval if and only if it joins one date row of the period. (b) A left join with at most one match keeps each left row once and supplies nothing that is read.}
{（a）由连续性，一个期间各月的键恰好是边界之间的整数区间，因此事实行的键落在区间内当且仅当它连接到该期间的一行日期。（b）至多一个匹配的左连接使每个左侧行恰好保留一次，且不提供任何被读取的内容。}
\end{proof}
\bi{The assumption of (a) is a condition of the new revision; that of (b) is a condition of the previous revision that the new revision no longer needs. Wherever the maintained conditions hold, a rule revision thus agrees with its previous revision on every period, and the sampled tests check the implementation and measure the cost. Maintenance checks an added assumption last; when it alone fails, every condition shared with the previous revision has just been confirmed, and \system\ invalidates the optimized revision and makes its previous revision current again.}
{（a）的前提是新版本的条件；（b）的前提是前一版本的条件，新版本不再需要它。因此，只要受维护的条件成立，规则版本就在每个期间上都与前一版本一致；抽样测试检查的是实现并测量代价。维护最后检查新增的前提；只有它不成立时，与前一版本共有的每个条件都刚刚得到确认，\system\ 使优化后的版本失效，并让前一版本重新成为当前版本。}

\subsection{\bt{Prediction tasks: shared feature definitions}{预测任务：共享的特征定义}}
\label{sec:features}
\bi{Predictive tasks reuse much of the same database knowledge as metric questions. Given an entity's seed rows and times, an agent writes SQL that aggregates each entity's earlier data; a fixed learner is then trained on those features. RelBench forecasting tasks follow this pattern~\cite{robinson2024relbench}. Much of the agent's work precedes its first feature validation, as it explores tables and probes values, so definitions learned in earlier sessions can provide an executable starting point.}
{预测任务与指标问题复用着许多相同的数据库知识。给定实体的种子行和时间，智能体编写 SQL，聚合每个实体此前的数据，再由固定学习器在这些特征上训练。RelBench 的预测任务采用这一形式~\cite{robinson2024relbench}。智能体的大量工作发生在第一次特征验证之前，用于查看表和探查取值，因此此前会话学得的定义可以提供可执行的起点。}

\bi{A \emph{feature definition} extends a metric's fact, measure, grain, verified many-to-one joins, and required filters with an entity, a look-back window, and a default value for seeds without rows. The entity is reached through a foreign key or verified join, optionally via an anchor such as a seed post's owner; the default matches the submitted SQL. Compilation produces one lateral aggregate per seed, bounding every contributing row by the seed time and window. These bounds prevent compiled features from reading data recorded after the seed time.}
{\emph{特征定义}在指标的事实表、度量、粒度、经验证的多对一连接和必需过滤之上，增加实体、回看窗口以及无匹配行时的默认值。实体经外键或验证过的连接到达，必要时经过种子帖子所有者这样的锚；默认值与提交的 SQL 一致。编译为每个种子生成一个 LATERAL 聚合，将每条参与计算的行限制在该种子的时间和窗口内，从而阻止编译出的特征读取种子时间之后记录的数据。}

\bi{A \emph{task program} retains the definitions, validation score, and complete SQL submitted by a successful session. The full SQL can express features beyond the structured form, such as self-joins. Consumers validate its predictive quality and audit temporal leakage, while the structured definitions receive the additional publication checks below.}
{\emph{任务方案}保留成功会话提交的定义、验证分数和完整 SQL。完整 SQL 还可以表达结构化形式之外的特征，例如自连接。使用者会验证其预测质量并审计时间泄漏，结构化定义则另外接受下面的发布检查。}
\bi{Feature learning follows the metric publication path. Once a session's features beat the entity-table reference, an extractor turns its SQL into definition drafts. On the current snapshot, each draft must have valid fields and a unique fact-grain key: columns must exist, expressions must read only the involved tables, and every join must be many-to-one. Its compiled SQL must also reproduce the submitted column value by value on 2,000 seeds. Recomputing on views that hide every row after seed time must leave those values unchanged.}
{特征学习沿用指标的发布流程。当会话的特征优于实体表参照时，提取器将 SQL 转为定义草稿。在当前快照上，每个草稿的字段必须合法、事实粒度键必须唯一：列需存在，表达式只能读取涉及的表，每个连接都需多对一。编译出的 SQL 还须在 2,000 个种子上逐值重现提交的列；在隐去种子时间之后所有行的视图上重算，结果也须保持不变。}

\bi{An agent starts with \code{find\_program} to retrieve its task's latest program. The response includes the full SQL after confirming that it runs on the current snapshot, together with SQL compiled from the program's still-valid definitions. To reuse work from other tasks, \code{find\_features} retrieves definitions published on the same entity by any agent, compiled for the current seed table; \code{compose\_features} compiles a chosen set. When data change, conditions are rechecked as in Section~\ref{sec:maintenance}. Feature maintenance revokes a definition whose condition fails, because these tasks provide no single reference value against which a repair could be verified.}
{智能体先通过 \code{find\_program} 检索本任务最近的方案。返回内容包括确认能在当前快照执行的完整 SQL，以及从方案中仍有效的定义编译得到的 SQL。若要复用其他任务的工作，\code{find\_features} 可检索任何智能体在同一实体上发布的定义，并为当前种子表编译；\code{compose\_features} 则编译选定的一组定义。数据变化后，条件按第~\ref{sec:maintenance} 节重新检查。特征维护会撤销条件失败的定义，因为这类任务没有单个参考值可以用于核实修复。}
